% ID: SEL-D0079-OA-E02
% Tanda S079.
\textbf{a)} Factorizamos la primera derivada y calculamos la segunda:
\[
 f'(x)=(3x-5)(x-1),\qquad f''(x)=6x-8.
\]
La primera derivada es positiva en $(-\infty,1)$ y $(5/3,\infty)$ y negativa en $(1,5/3)$. Por tanto, $f$ crece en los dos primeros intervalos y decrece en el último. Tiene un máximo relativo en $x=1$ y un mínimo relativo en $x=5/3$.

Como $f''<0$ para $x<4/3$, la gráfica se curva hacia abajo en $(-\infty,4/3)$. Como $f''>0$ para $x>4/3$, se curva hacia arriba en $(4/3,\infty)$. Hay un punto de inflexión de abscisa $4/3$.

\textbf{b)} La pendiente en el punto dado es $f'(1)=3-8+5=0$. Con $f(1)=1$, la ecuación punto-pendiente es $y-1=0(x-1)$, luego la tangente es $\boxed{y=1}$.
